\documentclass[11pt]{article}
\newcommand{\ListaTitulo}{Gabarito 10: Bernoulli, Binomial e Poisson}
% Preâmbulo compartilhado: MATF14, Listas de Exercícios 2026
% Usado por ListaNN.tex e GabaritoNN.tex via % Preâmbulo compartilhado: MATF14, Listas de Exercícios 2026
% Usado por ListaNN.tex e GabaritoNN.tex via % Preâmbulo compartilhado: MATF14, Listas de Exercícios 2026
% Usado por ListaNN.tex e GabaritoNN.tex via \input{preamble}

\usepackage[utf8]{inputenc}
\usepackage[T1]{fontenc}
\usepackage[brazil]{babel}
\usepackage{fullpage}
\usepackage{amsmath,amssymb}
\usepackage{enumitem}
\usepackage{xcolor}
\usepackage{fancyhdr}
\usepackage{titlesec}
\usepackage{listings}
\usepackage{hyperref}
\usepackage{booktabs}
\usepackage{graphicx}
\usepackage{tikz}

\definecolor{matf14azul}{HTML}{0D3B54}
\definecolor{matf14dourado}{HTML}{C9971E}

\titleformat{\section}{\normalfont\Large\bfseries\color{matf14azul}}{\thesection}{1em}{}
\titleformat{\subsection}{\normalfont\large\bfseries\color{matf14azul}}{\thesubsection}{1em}{}

\providecommand{\ListaTitulo}{Lista de Exercícios}

\setlength{\headheight}{14pt}
\pagestyle{fancy}
\fancyhf{}
\lhead{\small MATF14: Estatística Econômica I}
\rhead{\small \ListaTitulo}
\cfoot{\thepage}
\renewcommand{\headrulewidth}{0.4pt}

\lstset{
  basicstyle=\ttfamily\small,
  breaklines=true,
  frame=single,
  columns=fullflexible,
  backgroundcolor=\color{gray!8},
  commentstyle=\color{matf14dourado},
  keywordstyle=\color{matf14azul}\bfseries,
  language=R
}

\newcounter{questaocounter}
\newenvironment{questao}[1]{%
  \stepcounter{questaocounter}%
  \bigskip\noindent\textbf{Questão #1.}\ %
}{\par}

\newcommand{\cabecalho}[2]{%
  \begin{center}
    {\Large\bfseries\color{matf14azul} #1}\\[0.3em]
    {\large #2}
  \end{center}
  \vspace{0.5em}
  \noindent\rule{\linewidth}{0.6pt}
  \vspace{0.5em}
}

\newenvironment{solucao}{%
  \par\smallskip\noindent\textit{\color{matf14dourado!80!black}Solução.}\ %
}{\hfill$\blacksquare$\par}


\usepackage[utf8]{inputenc}
\usepackage[T1]{fontenc}
\usepackage[brazil]{babel}
\usepackage{fullpage}
\usepackage{amsmath,amssymb}
\usepackage{enumitem}
\usepackage{xcolor}
\usepackage{fancyhdr}
\usepackage{titlesec}
\usepackage{listings}
\usepackage{hyperref}
\usepackage{booktabs}
\usepackage{graphicx}
\usepackage{tikz}

\definecolor{matf14azul}{HTML}{0D3B54}
\definecolor{matf14dourado}{HTML}{C9971E}

\titleformat{\section}{\normalfont\Large\bfseries\color{matf14azul}}{\thesection}{1em}{}
\titleformat{\subsection}{\normalfont\large\bfseries\color{matf14azul}}{\thesubsection}{1em}{}

\providecommand{\ListaTitulo}{Lista de Exercícios}

\setlength{\headheight}{14pt}
\pagestyle{fancy}
\fancyhf{}
\lhead{\small MATF14: Estatística Econômica I}
\rhead{\small \ListaTitulo}
\cfoot{\thepage}
\renewcommand{\headrulewidth}{0.4pt}

\lstset{
  basicstyle=\ttfamily\small,
  breaklines=true,
  frame=single,
  columns=fullflexible,
  backgroundcolor=\color{gray!8},
  commentstyle=\color{matf14dourado},
  keywordstyle=\color{matf14azul}\bfseries,
  language=R
}

\newcounter{questaocounter}
\newenvironment{questao}[1]{%
  \stepcounter{questaocounter}%
  \bigskip\noindent\textbf{Questão #1.}\ %
}{\par}

\newcommand{\cabecalho}[2]{%
  \begin{center}
    {\Large\bfseries\color{matf14azul} #1}\\[0.3em]
    {\large #2}
  \end{center}
  \vspace{0.5em}
  \noindent\rule{\linewidth}{0.6pt}
  \vspace{0.5em}
}

\newenvironment{solucao}{%
  \par\smallskip\noindent\textit{\color{matf14dourado!80!black}Solução.}\ %
}{\hfill$\blacksquare$\par}


\usepackage[utf8]{inputenc}
\usepackage[T1]{fontenc}
\usepackage[brazil]{babel}
\usepackage{fullpage}
\usepackage{amsmath,amssymb}
\usepackage{enumitem}
\usepackage{xcolor}
\usepackage{fancyhdr}
\usepackage{titlesec}
\usepackage{listings}
\usepackage{hyperref}
\usepackage{booktabs}
\usepackage{graphicx}
\usepackage{tikz}

\definecolor{matf14azul}{HTML}{0D3B54}
\definecolor{matf14dourado}{HTML}{C9971E}

\titleformat{\section}{\normalfont\Large\bfseries\color{matf14azul}}{\thesection}{1em}{}
\titleformat{\subsection}{\normalfont\large\bfseries\color{matf14azul}}{\thesubsection}{1em}{}

\providecommand{\ListaTitulo}{Lista de Exercícios}

\setlength{\headheight}{14pt}
\pagestyle{fancy}
\fancyhf{}
\lhead{\small MATF14: Estatística Econômica I}
\rhead{\small \ListaTitulo}
\cfoot{\thepage}
\renewcommand{\headrulewidth}{0.4pt}

\lstset{
  basicstyle=\ttfamily\small,
  breaklines=true,
  frame=single,
  columns=fullflexible,
  backgroundcolor=\color{gray!8},
  commentstyle=\color{matf14dourado},
  keywordstyle=\color{matf14azul}\bfseries,
  language=R
}

\newcounter{questaocounter}
\newenvironment{questao}[1]{%
  \stepcounter{questaocounter}%
  \bigskip\noindent\textbf{Questão #1.}\ %
}{\par}

\newcommand{\cabecalho}[2]{%
  \begin{center}
    {\Large\bfseries\color{matf14azul} #1}\\[0.3em]
    {\large #2}
  \end{center}
  \vspace{0.5em}
  \noindent\rule{\linewidth}{0.6pt}
  \vspace{0.5em}
}

\newenvironment{solucao}{%
  \par\smallskip\noindent\textit{\color{matf14dourado!80!black}Solução.}\ %
}{\hfill$\blacksquare$\par}


\begin{document}

\cabecalho{Gabarito 10}{Distribuições de Bernoulli, Binomial e Poisson (Cap. 3, Aulas 20--21)}

\begin{questao}{1} \textbf{(Binomial: carteira de investidores)}
\begin{solucao}
$X\sim\text{Bin}(12;\,0{,}3)$.
\begin{enumerate}[label=(\alph*)]
  \item $P(X=4) = \binom{12}{4}(0{,}3)^4(0{,}7)^8 \approx 0{,}2311$.
  \item $P(X\le2) = \sum_{k=0}^2 \binom{12}{k}(0{,}3)^k(0{,}7)^{12-k} \approx 0{,}2528$.
  \item $E(X)=np=12\times0{,}3=3{,}6$. $\mathrm{Var}(X)=np(1-p)=12\times0{,}3\times0{,}7=2{,}52$;
    $dp(X)=\sqrt{2{,}52}\approx1{,}587$.
\end{enumerate}
\end{solucao}
\end{questao}

\begin{questao}{2} \textbf{(Binomial: controle de qualidade)}
\begin{solucao}
\begin{enumerate}[label=(\alph*)]
  \item $X\sim\text{Bin}(25;\,0{,}04)$. $P(X\le1) = P(X=0)+P(X=1) = (0{,}96)^{25} +
    25(0{,}04)(0{,}96)^{24} \approx 0{,}3604+0{,}3754 = 0{,}7358$.
  \item Com $p=0{,}10$: $X\sim\text{Bin}(25;\,0{,}10)$. $P(X\le1) = (0{,}90)^{25} +
    25(0{,}10)(0{,}90)^{24} \approx 0{,}0718+0{,}1994=0{,}2712$. A probabilidade de aceitação caiu
    bastante (de $\approx73{,}6\%$ para $\approx27{,}1\%$), o plano de amostragem ficou (como
    deveria) muito mais eficaz em \textbf{rejeitar} lotes de qualidade pior, mesmo mantendo a
    mesma regra de decisão.
\end{enumerate}
\end{solucao}
\end{questao}

\begin{questao}{3} \textbf{(Poisson: fila de atendimento)}
\begin{solucao}
\begin{enumerate}[label=(\alph*)]
  \item $X\sim\text{Poisson}(6)$. $P(X=4) = \dfrac{e^{-6}6^4}{4!} \approx 0{,}1339$.
  \item $P(X>8) = 1-P(X\le8) \approx 1-0{,}8472 = 0{,}1528$.
  \item Para meia hora, a taxa média é proporcional ao intervalo: $\lambda_{1/2h}=6/2=3$. Logo,
    $Y\sim\text{Poisson}(3)$ e $P(Y=2) = \dfrac{e^{-3}3^2}{2!} \approx 0{,}2240$.
\end{enumerate}
\end{solucao}
\end{questao}

\begin{questao}{4} \textbf{(Aproximação Binomial $\to$ Poisson: fraude rara)}
\begin{solucao}
\begin{enumerate}[label=(\alph*)]
  \item $X\sim\text{Bin}(5000;\,0{,}0006)$. $P(X=3) = \binom{5000}{3}(0{,}0006)^3(0{,}9994)^{4997}
    \approx 0{,}2222$.
  \item $\lambda=np=5000\times0{,}0006=3$. $P(X=3)\approx \dfrac{e^{-3}3^3}{3!}\approx 0{,}2240$.
  \item Os dois valores são muito próximos ($0{,}2222$ vs. $0{,}2240$). A aproximação funciona bem
    porque $n=5000$ é grande e $p=0{,}0006$ é pequeno, exatamente as condições vistas em aula
    para a aproximação de Poisson à Binomial, com $\lambda=np$ moderado (aqui, $\lambda=3$).
\end{enumerate}
\end{solucao}
\end{questao}

\begin{questao}{5} \textbf{(Dedução: esperança da Binomial a partir da soma de Bernoullis)}
\begin{solucao}
\begin{enumerate}[label=(\alph*)]
  \item
  $$
  E(X) = E(X_1+\cdots+X_n) = E(X_1)+\cdots+E(X_n) = \underbrace{p+p+\cdots+p}_{n \text{ vezes}} = np.
  $$
  \item Como as $X_i$ são independentes,
  $$
  \mathrm{Var}(X) = \mathrm{Var}(X_1+\cdots+X_n) = \mathrm{Var}(X_1)+\cdots+\mathrm{Var}(X_n)
  = \underbrace{p(1-p)+\cdots+p(1-p)}_{n \text{ vezes}} = np(1-p).
  $$
  \item A independência é usada explicitamente na passagem
    $\mathrm{Var}(X_1+\cdots+X_n)=\sum_i\mathrm{Var}(X_i)$ do item (b), essa igualdade só vale
    porque as $X_i$ são independentes (o caso geral teria termos adicionais de covariância entre
    pares $X_i,X_j$). Já no item (a), a igualdade $E(X_1+\cdots+X_n)=\sum_i E(X_i)$ é usada sem
    exigir independência, essa propriedade da esperança vale \textbf{sempre}, dependentes ou não.
    É por isso que $E(X)=np$ vale mesmo se, hipoteticamente, as tentativas não fossem
    independentes, mas $\mathrm{Var}(X)=np(1-p)$ deixaria de valer nesse cenário.
\end{enumerate}
\end{solucao}
\end{questao}

\end{document}
